Everything is computed in Gauge Theory Gravity (Definition \claimref{obj.synthesis.gtg-flat-spacetime-gauge-fields}; Fact \claimref{synthesis.imported-gtg-dictionary}, clauses (i), (iv) and (v)), on the flat Euclidean background $\Cl(4,0)$ for (i) to (iii) and on $\Cl(1,3)$ for (iv) and (v); the closed manifolds of (ii) lie beyond one flat chart, and only their pointwise densities are computed on a chart. \emph{Notation.} For curvature operators $X, Y$ (linear maps on bivectors, such as the field strength $\mathcal{R}(B) = -R(\hb(B))$) put
\[
\mathcal{Q}(X, Y) \;=\; 2\,\langle X(a\wedge b)\,I^{-1}\,Y(I\,\partial_{b}\wedge\partial_{a})\rangle ,
\]
symmetric and bilinear; $\epsilon_{abcd}X^{ab}\wedge Y^{cd} = \mathcal{Q}(X, Y)\,\mathrm{vol}$ for the corresponding two-forms in Euclidean signature, and $\mathcal{Q}(X) := \mathcal{Q}(X, X)$ is the Gauss-Bonnet integrand $E(X)$ (dictionary row 13 in the $I^{-1}$ form, which is correct in both signatures). Let $D$ be the curvature operator of the unit four-sphere; in the store's sign on $\Cl(4,0)$ it is $D(B) = -B$ (the check solves the round $\Sph^{4}$ as a gauge field and finds $\mathcal{R}(a) = -3a$), and put $s(X) := \tfrac12\mathcal{Q}(X, D)$, which is the scalar trace in the convention positive on spheres ($s(D) = 12$; $s = -\partial_{a}\cdot\mathcal{R}(a)$ here). Then $\mathcal{Q}(D) = 24$. With $F = \mathcal{R} - \ell^{-2}D$ the Lorentz block of the Euclidean de Sitter-branch curvature (Fact \claimref{dark-energy.imported-macdowell-mansouri-theorem-stelle-west}; the sign is the one for which the round $\Sph^{4}$ of radius $\ell$ has $F = 0$; the overall sign of $F$ does not enter $\mathcal{Q}(F)$), bilinearity gives
\begin{equation}
\mathcal{Q}(F) \;=\; E(\mathcal{R}) - 4\ell^{-2}s(\mathcal{R}) + 24\,\ell^{-4}. \tag{1}
\end{equation}
The check computes $\mathcal{Q}$ in $\Cl(4,0)$ from the field strengths of three solved gauge fields and finds $\mathcal{Q}(D) = 24$, $\mathcal{Q}(\mathcal{R}, D) = 2s(\mathcal{R})$ and $\mathcal{Q}(\mathcal{R}) = E(\mathcal{R})$.

\emph{(i)} By (1), $S = k\ell^{2}\int\mathcal{Q}(F)\,\mathrm{vol} = k\ell^{2}\int E - 4k\int s + 24k\ell^{-2}\int\mathrm{vol}$, the three terms of the Fact's decomposition. Matching $-4k\int s$ to $-\tfrac{1}{16\pi G}\int s$ gives $k = 1/(64\pi G)$, and then $24k\ell^{-2} = 3/(8\pi G\ell^{2})$, which equals $2\Lambda/(16\pi G)$ exactly when $\Lambda = 3/\ell^{2}$, the value the Fact's field equations carry; the Gauss-Bonnet coefficient is $k\ell^{2} = \ell^{2}/(64\pi G) = 3/(64\pi G\Lambda)$.

\emph{(ii)} Let the field be Einstein, $\mathcal{R}(a) = -\Lambda a$ on $\Cl(4,0)$ in the store's sign ($\Ric = \Lambda g$ in the convention positive on spheres). The decomposition of the curvature operator into its Weyl and scalar parts (the trace-free Ricci part vanishing) gives $\mathcal{R} = \mathcal{W} + \tfrac{\Lambda}{3}D$, so $F = \mathcal{W} + (\tfrac{\Lambda}{3} - \ell^{-2})D = \mathcal{W}$. The Weyl operator commutes with the duality, $\mathcal{W}(IB) = I\mathcal{W}(B)$ (dictionary row 5), so
\[
\mathcal{Q}(F) = 2\langle\mathcal{W}(a\wedge b)\,I^{-1}I\,\mathcal{W}(\partial_{b}\wedge\partial_{a})\rangle = 2\langle\mathcal{W}(a\wedge b)\,\mathcal{W}(\partial_{b}\wedge\partial_{a})\rangle = |W|^{2} .
\]
Integrating, $I_{\mathrm{MM}}[M] = k\ell^{2}\int_{M}|W|^{2}\mathrm{vol} \ge 0$, with equality exactly when $\mathcal{W} \equiv 0$, which in dimension four is local conformal flatness. For the second form: $E(\mathcal{R}) = \mathcal{Q}(\mathcal{W}) + \tfrac{2\Lambda}{3}\mathcal{Q}(\mathcal{W}, D) + \tfrac{\Lambda^{2}}{9}\mathcal{Q}(D) = |W|^{2} + \tfrac83\Lambda^{2}$, since $\mathcal{Q}(\mathcal{W}, D) = 2s(\mathcal{W}) = 0$ ($\mathcal{W}$ is trace-free), and the Chern-Gauss-Bonnet theorem $\int_{M}E = 32\pi^{2}\chi(M)$ (Fact \claimref{synthesis.imported-chern-gauss-bonnet}(b)) gives
\[
I_{\mathrm{EH}} + \frac{\chi}{2}\pi N^{2} \;=\; -\frac{\Lambda\,\mathrm{Vol}}{8\pi G} + k\ell^{2}\int\Bigl(|W|^{2} + \tfrac83\Lambda^{2}\Bigr) \;=\; k\ell^{2}\int|W|^{2} + \Bigl(-\frac{\Lambda}{8\pi G} + \frac{3}{64\pi G\Lambda}\cdot\frac{8\Lambda^{2}}{3}\Bigr)\mathrm{Vol} \;=\; I_{\mathrm{MM}},
\]
using $I_{\mathrm{EH}} = -\tfrac{1}{16\pi G}\int(s - 2\Lambda) = -\Lambda\mathrm{Vol}/(8\pi G)$ on an Einstein space and $k\ell^{2}\cdot 32\pi^{2}\chi = \pi\ell^{2}\chi/(2G) = \tfrac{\chi}{2}\pi N^{2}$ with $N^{2} = \ell^{2}/G$ (Proposition \claimref{synthesis.de-sitter-entropy-pi-n-squared}: $N = M_{P}/H$, $H^{2} = \Lambda/3$). In units of $\pi N^{2}$, $I_{\mathrm{EH}} = -\Lambda^{2}\mathrm{Vol}/(24\pi^{2})$ and $I_{\mathrm{MM}} = \mathrm{Vol}\,|W|^{2}/(64\pi^{2})$, both rational once $\mathrm{Vol}/\pi^{2}$ is. The instantons are homogeneous, so $E$ and $|W|^{2}$ are constant on each and the Chern-Gauss-Bonnet theorem gives $\mathrm{Vol} = 32\pi^{2}\chi/E$; then $I_{\mathrm{MM}}/\pi N^{2} = \chi|W|^{2}/(2E)$ and $I_{\mathrm{EH}}/\pi N^{2} = -4\chi\Lambda^{2}/(3E)$, all scale invariant. The check builds each as a gauge field on $\Cl(4,0)$ and solves its rotation-gauge field from the torsion equation: the round $\Sph^{4}$ in the stereographic gauge $\hb^{-1}(e_{i}) = \tfrac{2}{1 + Q(x)}\,e_{i}$; $\Sph^{2}\times\Sph^{2}$ with one stereographic factor on $(e_{1}, e_{2})$ and one on $(e_{3}, e_{4})$; and $\mathbb{CP}^{2}$ (Fubini-Study) in the gauge $f^{0} = \nabla r/(1 + r^{2})$, $f^{1,2} = r\sigma_{1,2}/(2\sqrt{1 + r^{2}})$, $f^{3} = r\sigma_{3}/(2(1 + r^{2}))$ with the left-invariant fields $\sigma_{l}$ of $SU(2)$ as in the Taub-NUT result, at a point where every value is rational. From the field strengths it finds the three fields Einstein ($\Lambda = 3, 1, 6$ in these normalisations), $\mathcal{Q}(F) = |W|^{2} = 0, \tfrac{16}{3}, 96$, $E = 24, 8, 192$ with $E = |W|^{2} + \tfrac83\Lambda^{2}$, and so, with $\chi = 2, 4, 3$, $I_{\mathrm{MM}}/\pi N^{2} = 0$, $\tfrac43$, $\tfrac34$ and $I_{\mathrm{EH}}/\pi N^{2} = -1$, $-\tfrac23$, $-\tfrac34$; the volumes this implies, $\tfrac83\pi^{2}$, $16\pi^{2}$ and $\tfrac12\pi^{2}$, are the standard ones (cited from memory). The pointwise identity $\mathcal{Q}(\mathcal{R} - \tfrac{\Lambda}{3}D) = |W|^{2}$ is also checked, in components, on a generic rational Weyl tensor (the Weyl part of a combination of Kulkarni-Nomizu products) at three values of $\Lambda$, both signs included.

\emph{(iii)} The promotion (Fact \claimref{dark-energy.imported-henneaux-teitelboim-promotion}) replaces the fixed number $\ell^{2}$ by a field $\lambda$ and adds $-\int\lambda\,d\mathcal{T}$ with $\mathcal{T}$ a trivector field and $d\mathcal{T} = \nabla\wedge\mathcal{T}$ (a three-form in components); varying $\mathcal{T}$ gives $\nabla\lambda = 0$, which is the equation the norm mode of Proposition \claimref{dark-energy.under-action-written-norm-parameter-laws} lacked, and varying $\lambda$ gives $d\mathcal{T} = \partial L/\partial\lambda$, so $\mathcal{T}$ advances between two slices by the integral of $\partial L/\partial\lambda$ over the region between them. At fixed connection and frame, by (1) with $u = \ell^{2}$, $L = k\bigl(uE - 4s + 24u^{-1}\bigr)\mathrm{vol}$ and
\[
\frac{\partial L}{\partial u} \;=\; k\bigl(E - 24u^{-2}\bigr)\mathrm{vol} \;=\; k\bigl(\mathcal{Q}(F) + 4u^{-1}\,s(F)\bigr)\mathrm{vol},
\]
the second form by (1) and $s(F) = s(R) - 12u^{-1}$ (the identity is checked in exact arithmetic at two values of $u$ on a curvature tensor that is not Einstein). The frame equation of the Fact is, on $\Cl(4,0)$, the trivector-valued $2\,\partial_{b}\wedge F(I(a\wedge b)) = 2(\mathcal{G}(a) - \Lambda a)I = 0$ ($\epsilon_{abcd}\,e^{b}\wedge F^{cd} = 0$ in components), by $\partial_{b}\wedge\mathcal{R}(I(a\wedge b)) = \mathcal{G}(a)I$, $\partial_{b}\wedge(I(a\wedge b)) = -3aI$ and $F = \mathcal{R} - \ell^{-2}D = \mathcal{R} + \ell^{-2}\,\mathrm{id}$ (on $\Cl(1,3)$, where the de Sitter operator is $+\tfrac{\Lambda}{3}\,\mathrm{id}$, the same computation gives $2(\mathcal{G}(a) + \Lambda a)I$); the check evaluates the trivector on the three instantons and finds it zero. Traced with the frame it gives $\partial_{a}\cdot\mathcal{G}(a) - 4\Lambda = s(\mathcal{R}) - 4\Lambda = 0$, so $s(\mathcal{R}) = 12u^{-1}$ and $\mathcal{Q}(F, D) = 2s(F) = 0$ on shell, so the density is $k\,\mathcal{Q}(F) = k|W|^{2}$ by (ii), torsion-free solutions being Einstein with $\Lambda = 3/\ell^{2}$. If instead a rescaled frame $\hat e = e/f(u)$ is held fixed, the derivative changes by $\int (\delta S_{\mathrm{tot}}/\delta e^{a})\wedge f'(u)\hat e^{a}$, with $S_{\mathrm{tot}}$ the whole action (the matter action included when there is one), a multiple of the trace of the full frame equation, which vanishes on shell. The density $k(E - 24u^{-2})$ itself is obtained at fixed frame and connection and uses no field equation; only its reduction to $k\,\mathcal{Q}(F)$, and then to $k|W|^{2}$, uses the vacuum equations, through $s(F) = 0$ and $F = W$. For any other promoted variable $g(u)$ the term $-\int g\,d\mathcal{T}_{g}$ gives $d\mathcal{T}_{g} = \partial L/\partial g = (\partial L/\partial u)/g'(u)$, so its conjugate advances by the same integral divided by $g'(u)$; for $\Lambda = 3/u$ this is $(du/d\Lambda) = -3/\Lambda^{2}$ times it, and for the norm $Q(I) = \Lambda/3 = 1/u$ it is $-u^{2}$ times it. No choice is singled out; on each solution the conjugate times are proportional. In Lorentzian signature the same algebra holds with $\epsilon_{abcd}\epsilon^{efgh} = -\delta^{efgh}_{abcd}$, so the density is a multiple of $W_{abcd}W^{abcd}\sqrt{|g|}$.

\emph{(iv)} On a vacuum solution the density is $k|W|^{2}$ by (iii), so it vanishes on regions where the Weyl part of the field strength vanishes. A vacuum solution is Einstein with $\Lambda = 3/\ell^{2}$, and Einstein with $W = 0$ means the whole curvature is the de Sitter operator $\mathcal{R}(B) = \tfrac{\Lambda}{3}B$ of the store's sign on $\Cl(1,3)$ (not $\tfrac{\Lambda}{3}D$: $D = -\mathrm{id}$ is the Euclidean sphere's operator, and $-\tfrac{\Lambda}{3}\,\mathrm{id}$ on $\Cl(1,3)$ is anti-de Sitter), constant curvature, so such a region is locally de Sitter; the de Sitter slicings of FLRW form are among them. (Every FLRW field has vanishing Weyl part, its derived metric being conformally flat: in the comoving orthonormal frame its curvature has $R_{0i0j} = A\,\delta_{ij}$ and $R_{ijkl} = B\,D_{ijkl}$, so the electric and magnetic parts of $W$ are $SO(3)$-invariant symmetric trace-free tensors and vanish; the earlier check computes the Lorentzian Weyl tensor at three rational points $(a, \dot a, \ddot a, k)$ and finds it zero, and the GTG check finds the Weyl operator of the flat FLRW gauge field $\hb^{-1}(e_{i}) = a\,e_{i}$ zero at five instants. But an FLRW metric that is not de Sitter is not a vacuum solution, and on it the density is not $k|W|^{2}$: see (v).) On the leading Kasner term of Proposition \claimref{dark-energy.frozen-point-datum-kasner-wall} the metric is Ricci-flat, so $W$ is the whole curvature and $|W|^{2}$ is the Kretschmann scalar $\bigl(4\sum_{i}p_{i}^{2}(p_{i} - 1)^{2} + 4\sum_{i<j}p_{i}^{2}p_{j}^{2}\bigr)\tau^{-4}$, which for $(\tfrac23, \tfrac23, -\tfrac13)$ is $\tfrac{64}{27}\tau^{-4}$, as that Proposition records (the GTG check computes $2\langle\mathcal{W}(a\wedge b)\mathcal{W}(\partial_{b}\wedge\partial_{a})\rangle$ from the Kasner gauge field $\hb^{-1}(e_{i}) = \tau^{p_{i}}e_{i}$ at $\tau = 1, 2, \tfrac13$ and finds it Ricci-flat with this value); with $\det(\hb)^{-1} = \tau^{p_{1}+p_{2}+p_{3}} = \tau$ the density is $\tfrac{64}{27}k\,\tau^{-3}$ and its integral from $\tau$ to a fixed $\tau_{0}$ grows as $\tau^{-2}$. The $\Lambda$-corrections are of relative order $\Lambda\tau^{2}$ by that Proposition and do not change the divergence. Under the Fact's identity $F = W$ on Einstein solutions this is exact, not only at leading order, wherever the solution is Einstein with the promoted $\Lambda$.

\emph{(v)} Under (Mat) the matter action does not depend on $u$ at fixed frame, connection and matter fields, so the variation in the multiplier gives the density $k(E - 24u^{-2})\mathrm{vol}$ of (iii) unchanged. The field equations of the functional plus the matter action are the Fact's equations with the matter source, Einstein's equations with $\Lambda = 3/\ell^{2}$ and stress tensor $T$ (the Gauss-Bonnet term contributes nothing to the Fact's Euler-Lagrange equations, which are linear in the source). Frame normalisation: in Lorentzian signature the functional is $-k(uE - 4s + 24u^{-1})\sqrt{|g|}$ plus the matter Lagrangian, the sign fixed by matching $4ks = s/(16\pi G)$ to the Einstein-Hilbert term; at fixed frame its $u$-derivative is $-k(E - 24u^{-2})\sqrt{|g|}$; holding the unit frame, $g = u\hat g$, the gravitational part gives $-k\,\mathcal{Q}(F)\sqrt{|g|}$ and the matter part gives $t\sqrt{|g|}/(2u)$, from $T^{ab} = (2/\sqrt{|g|})\,\partial(\sqrt{|g|}L_{m})/\partial g_{ab}$; the difference is $-(\sqrt{|g|}/u)\bigl(4k\,s(F) + t/2\bigr)$, which vanishes because $s(F) = s - 12u^{-1} = -8\pi G\,t$ is the trace of Einstein's equations. The check verifies the equality, and that the gravitational part alone differs from it whenever $t \neq 0$, at three FLRW points (units $8\pi G = 1$, $k = 1/8$).

In four dimensions the curvature decomposes into Weyl, trace-free Ricci and scalar parts, the Weyl tensor plus the Kulkarni-Nomizu products of $\Ric^{\mathrm{tf}}$ with $g$ and of $g$ with itself, the three parts orthogonal, and $E = |W|^{2} - 2|\Ric^{\mathrm{tf}}|^{2} + \tfrac16 s^{2}$, with norms taken with the Lorentzian metric (the check verifies this on a generic Lorentzian algebraic curvature tensor). Einstein's equations give, from their trace and trace-free parts, $s = 4\Lambda - 8\pi G\,t$ and $\Ric^{\mathrm{tf}} = 8\pi G\,T^{\mathrm{tf}}$, and $\tfrac16 s^{2} - \tfrac83\Lambda^{2} = -\tfrac43\Lambda(8\pi G\,t) + \tfrac16(8\pi G\,t)^{2}$, which is the first display; the check verifies it on the generic tensor with $T$ defined by Einstein's equations at three values of $\Lambda$. For flat FLRW, in the comoving orthonormal frame, $R_{0i0j} = -q\,\delta_{ij}$ and $R_{ijkl} = h\,D_{ijkl}$ with $h = H^{2}$ and $q = \ddot a/a$; then $|R|^{2} = 12(q^{2} + h^{2})$, $\Ric_{00} = -3q$, $\Ric_{ii} = q + 2h$, $s = 6(q + h)$, and $E = 36(q + h)^{2} - 4\bigl(9q^{2} + 3(q + 2h)^{2}\bigr) + 12(q^{2} + h^{2}) = 24hq$. Einstein's equations with $T = \mathrm{diag}(\rho, p, p, p)$ are the Friedmann equations $h = (\Lambda + \tilde\rho)/3$, $q = \Lambda/3 - (\tilde\rho + 3\tilde p)/6$, and
\[
24hq - \tfrac83\Lambda^{2} = \tfrac{4}{3}\bigl[(\Lambda + \tilde\rho)(2\Lambda - \tilde\rho - 3\tilde p) - 2\Lambda^{2}\bigr] = \tfrac43\bigl[\Lambda(\tilde\rho - 3\tilde p) - \tilde\rho(\tilde\rho + 3\tilde p)\bigr].
\]
In GTG the same values come from the flat FLRW gauge field with $\dot a/a = H$, $\ddot a/a = q$: the GTG check computes its field strength at five rational $(\Lambda, \tilde\rho, \tilde p)$ with $H^{2} = (\Lambda + \tilde\rho)/3$ and $q$ from the second Friedmann equation, finds the store-sign field equation $\mathcal{G}(a) + \Lambda a = -\tilde T(a)$ with $\tilde T(a) = (\tilde\rho + \tilde p)(a\cdot e_{0})e_{0} - \tilde p\,a$, $E = 24hq = \mathcal{Q}(\mathcal{R})$ and the display. The earlier check computes $E$ from the tensor, verifies both Einstein components from the same tensor, and verifies the display on the $3 \times 3 \times 3$ grid $\{\tfrac12, 3, -2\}^{3}$ of $(\Lambda, \tilde\rho, \tilde p)$; both sides have degree at most two in each variable, so agreement on three values per variable is the polynomial identity. Radiation ($\tilde p = \tilde\rho/3$) gives $-\tfrac83\tilde\rho^{2}$, and dust gives $\tfrac43\tilde\rho(\Lambda - \tilde\rho)$, positive for $0 < \tilde\rho < \Lambda$ and negative for $\tilde\rho > \Lambda$. Finally $(3k/\Lambda^{2})\bigl(-\tfrac43\Lambda\cdot 8\pi G\,t\bigr) = -t/(2\Lambda)$ with $k = 1/(64\pi G)$, and the remaining terms carry a further factor $G$; for dust $t = -\rho$, and $\int\rho\,\sqrt{|g|}\,d^{4}x$ is the sum over the dust's worldlines of rest mass times proper time.

\emph{Reading of (v) (argued, not part of the statement).} In a radiation era $\tilde\rho \propto a^{-4}$ and $H \propto a^{-2}$, so per comoving volume the advance between scale factors $a_{1} < a_{2}$ is proportional to $\int \tilde\rho^{2} a^{3}\,dt = \int \tilde\rho^{2}a^{2}H^{-1}\,da \propto \int a^{-4}\,da$, that is to $a_{1}^{-3} - a_{2}^{-3}$: the record is dominated by the earliest epoch. For dust after $8\pi G\rho = \Lambda$ the density changes sign, and as $\tilde\rho \to 0$ the advance per comoving volume tends to $\tfrac43 k\Lambda\,\tilde\rho\, a^{3}$ per unit cosmic time, a constant, since $\rho a^{3}$ is. The advance is therefore not monotone over a history with radiation and dust (negative before $8\pi G\rho_{m} \approx \Lambda$, positive after), so the conjugate variable is a record of the history and not a time coordinate that orders it. Against the Henneaux-Teitelboim clock of the Einstein-Hilbert action, which counts four-volume over $8\pi G$, the first-order term is that clock times $8\pi G\rho/(2\Lambda) = \rho/(2\rho_{\Lambda})$ with $\rho_{\Lambda} = \Lambda/8\pi G$.

\emph{The reading under (H).} If the weight of a value of the promoted $\Lambda$ is $e^{-I}$ summed over closed Einstein instantons, (ii) gives $e^{-I_{\mathrm{MM}}[\Sph^{4}]} = 1$ at every $\Lambda$ and $e^{-I_{\mathrm{MM}}[M]} = \exp\bigl(-\tfrac{3}{64\pi G\Lambda}\int|W|^{2}\bigr) < 1$ otherwise, with $\int|W|^{2}$ scale invariant, so these factors decrease as $\Lambda$ decreases. The Einstein-Hilbert weight of $\Sph^{4}$, $e^{\pi N^{2}}$, which is the factor used to argue for a vanishing cosmological constant (Baum, Hawking, Coleman; cited from memory), is multiplied here by $e^{-\pi N^{2}}$ from the Gauss-Bonnet term. So under (H) no factor of the action favours small $\Lambda$, and the measure is not supplied by the action. (H) itself, and the Euclidean continuation of the functional, are hypotheses of the reading, not results; the Euclidean value $I_{\mathrm{MM}}[\Sph^{4}] = 0$ is the Euclidean counterpart of the remark in Proposition \claimref{synthesis.de-sitter-entropy-pi-n-squared} that the functional as written gives the horizon zero entropy.


\emph{Literature.} The identity $I_{\mathrm{MM}} = k\ell^{2}\int|W|^{2}$ on Einstein spaces is known for the Euler-regularised action of anti-de Sitter gravity, where it expresses the on-shell equivalence of that action with conformal gravity (O.~Miskovic and R.~Olea, ``Topological regularization and self-duality in four-dimensional anti-de Sitter gravity'', arXiv:0902.2082; Fact \claimref{synthesis.imported-miskovic-olea-identity}); the computation above is the de Sitter (Euclidean, positive $\Lambda$) form of the same algebra.

\emph{In components} (a labelled reading). In an orthonormal frame a curvature-type two-form is $X^{ab} = \tfrac12 X^{ab}{}_{ef}\,e^{e}\wedge e^{f}$, and in Euclidean signature $\epsilon_{abcd}\epsilon^{efgh} = \delta^{efgh}_{abcd}$ gives $\mathcal{Q}(X, Y) = \tfrac14\,\delta^{efgh}_{abcd}\,X^{ab}{}_{ef}\,Y^{cd}{}_{gh}$, with $E(X) = s(X)^{2} - 4|\Ric(X)|^{2} + |X|^{2}$, $D_{abcd} = \delta_{ac}\delta_{bd} - \delta_{ad}\delta_{bc}$ and $\delta^{efgh}_{abgh} = 2\delta^{ef}_{ab}$; in Lorentzian signature $\epsilon_{abcd}\epsilon^{efgh} = -\delta^{efgh}_{abcd}$, and the density is a multiple of $W_{abcd}W^{abcd}\sqrt{|g|}$. The earlier checks compute these contractions on the curvature tensors of the three instantons, with the Levi-Civita symbol read off the Clifford product of the frame vectors, and give the same values.
